% sec6_5.tex — Section 6.5: Asymptotics for Sample Means of Serially Correlated Processes

This section studies the asymptotic properties of the sample mean
\[
  \bar{y} \equiv \frac{1}{n} \sum_{t=1}^{n} y_t
\]
for serially correlated processes. (It should be kept in mind that $\bar{y}$ depends on
$n$, although the notation does not make it explicit.) For the consistency of the
sample mean, we already have a sufficient condition in the Ergodic Theorem of
Chapter~2. The first proposition of this section provides another sufficient condition
in the form of restrictions on covariance-stationary processes. We also provide
two CLTs for serially correlated processes, one for linear processes and the other
for ergodic stationary processes. Chapter~2 includes a CLT for ergodic stationary
processes, but it rules out serial correlation because the process is assumed to be a
martingale difference sequence. The CLT for ergodic stationary processes of this
section generalizes it by allowing for serial correlation.

\subsection*{LLN for Covariance-Stationary Processes}

\begin{proposition}[LLN for covariance-stationary processes with vanishing autocovariances]
\label{prop:6.8}
Let $\{y_t\}$ be covariance-stationary with mean $\mu$ and $\{\gamma_j\}$ be the autocovariances
of $\{y_t\}$. Then
\begin{enumerate}[label=(\alph*)]
\item $\bar{y} \to_{\mathrm{m.s.}} \mu$ as $n \to \infty$ if $\lim_{j \to \infty} \gamma_j = 0$.

\item $\displaystyle \lim_{n \to \infty} \Var(\sqrt{n}\,\bar{y})
  = \sum_{j=-\infty}^{\infty} \gamma_j < \infty$ \quad
  if $\{\gamma_j\}$ is summable.\footnote{Anderson (1971, Lemma~8.3.1, p.~460). Absolute summability of $\{\gamma_j\}$ is not needed.}
\end{enumerate}
\end{proposition}

Since a mean square convergence implies a convergence in probability, part~(a)
shows that a very mild condition on covariance-stationary processes suffices for
the sample mean to be consistent for $\mu$. To prove part~(a), by Chebychev's LLN
(see Section~2.1), it suffices to show that $\lim_{n \to \infty} \Var(\bar{y}) = 0$, which is fairly easy
(Analytical Exercise~9) once the expression for $\Var(\bar{y})$ is derived. Unlike in the case
without serial correlation in $\{y_t\}$, calculating the variance of $\bar{y}$ is more cumbersome
because covariances between two different terms have to be taken into account:
\begin{align}
  n \cdot \Var(\bar{y}) &= \Var(\sqrt{n}\,\bar{y}) \notag \\
  &= \frac{1}{n} [\Cov(y_1, y_1 + \cdots + y_n) + \cdots + \Cov(y_n, y_1 + \cdots + y_n)] \notag \\
  &= \frac{1}{n} [(\gamma_0 + \gamma_1 + \cdots + \gamma_{n-1}) + (\gamma_1 + \gamma_0 + \gamma_1 + \cdots + \gamma_{n-2}) \notag \\
  &\qquad + \cdots + (\gamma_{n-1} + \gamma_{n-2} + \cdots + \gamma_1 + \gamma_0)] \notag \\
  &= \frac{1}{n} [n\gamma_0 + 2(n-1)\gamma_1 + \cdots + 2(n-j)\gamma_j + \cdots + 2\gamma_{n-1}] \notag \\
  &= \gamma_0 + 2 \sum_{j=1}^{n-1} \Bigl(1 - \frac{j}{n}\Bigr) \gamma_j.
  \label{eq:6.5.2}
\end{align}
(If you have difficulty deriving this, set $n$ to, say, 3 to verify.) For each fixed $j \geq 1$,
the coefficient of $\gamma_j$ in \eqref{eq:6.5.2} goes to 1 as $n \to \infty$, so that all the autocovariances
eventually enter the sum with a coefficient of one. This is not enough to show the
desired result \eqref{eq:6.5.1}. Filling in the rest of the proof of (b) is Analytical Exercise~10.

For a covariance-stationary process $\{y_t\}$, we define the \textbf{long-run variance} to
be the limit as $n \to \infty$ of $\Var(\sqrt{n}\,\bar{y})$ (if it exists). So part~(b) of the proposition
says that the long-run variance equals $\sum_{j=-\infty}^{\infty} \gamma_j$, which in turn equals the value of
the autocovariance-generating function $g_Y(z)$ at $z = 1$ (see \eqref{eq:6.2.15}) or $2\pi$ times
the spectrum at frequency zero (see \eqref{eq:6.2.17}).

\subsection*{Two Central Limit Theorems}

We present two CLTs to cover serial correlation. The first is a generalization of
Lindeberg--Levy.

\begin{proposition}[CLT for MA($\infty$)]
\label{prop:6.9}
Let $y_t = \mu + \sum_{j=0}^{\infty} \psi_j \varepsilon_{t-j}$ where $\{\varepsilon_t\}$ is
independent white noise and $\sum_{j=0}^{n} |\psi_j| < \infty$. Then
\begin{equation}
  \sqrt{n}\,(\bar{y} - \mu) \dto N\Bigl(0,\; \sum_{j=-\infty}^{\infty} \gamma_j\Bigr).
  \label{eq:6.5.3}
\end{equation}
\end{proposition}

We will not prove this result (see Anderson, 1971, Theorem~7.7.8 or Brockwell
and Davis, 1991, Theorem~7.1.2 for a proof), but we should not be surprised to
see that $\avar(\bar{y})$ (the variance of the limiting distribution of $\sqrt{n}\,(\bar{y} - \mu)$) equals
the long-run variance $\sum_{j=-\infty}^{\infty} \gamma_j$. We know from Lemma~2.1 that ``$x_n \dto x$,''
``$\Var(x_n) \to c < \infty$'' $\Rightarrow$ ``$\Var(x) = c$.'' Here, set $x_n = \sqrt{n}\,(\bar{y} - \mu)$. By \eqref{eq:6.5.1},
$\Var(x_n)$ converges to a finite limit $\sum \gamma_j$. So the variance of the limiting distribution
of $\sqrt{n}\,(\bar{y} - \mu)$ is $\sum \gamma_j$.

By Proposition~6.1(d), the process in Proposition~6.9 is (stationary and) ergodic.
Remember that ergodicity limits serial correlation by requiring that two random
variables sufficiently far apart in the sequence be almost independent. But ergodic
stationarity alone is not sufficient to ensure the asymptotic normality of $\bar{y}$. A natural
question arises: is there a further restriction on ergodicity, short of requiring the
process to be linear as in Proposition~6.9, that delivers asymptotic normality? The
restriction that has become increasingly popular is what this book calls \textbf{Gordin's
condition} for ergodic stationary processes. It has three parts. The first two parts
follow:
\begin{enumerate}[label=(\alph*)]
\item $\E(y_t^2) < \infty$. (This is a restriction on [strictly] stationary processes because,
strictly speaking, a stationary process might not have finite second moments.)

\item $\E(y_t \mid y_{t-j}, y_{t-j-1}, \ldots) \to_{\mathrm{m.s.}} 0$ as $j \to \infty$. (Since $\{y_t\}$ is stationary, this
condition is equivalent to $\E(y_0 \mid y_{-j}, y_{-j-1}, \ldots) \to_{\mathrm{m.s.}} 0$ as $j \to \infty$. Because
$\E(y_t \mid y_{t-j}, y_{t-j-1}, \ldots)$ is a random variable, the convergence cannot be the
usual convergence for real numbers.)
\end{enumerate}

This condition implies that the unconditional mean is zero, $\E(y_t) = 0$.\footnote{See Lemma~5.14 of White (1984) for a proof.}
As the forecast (the conditional expectation) is based on less and less information, it should
approach the forecast based on no information, i.e., the unconditional expectation.
To prepare for the third part, let $I_t \equiv (y_t, y_{t-1}, y_{t-2}, \ldots)$ and write $y_t$ as
\begin{align*}
  y_t &= y_t - [\E(y_t \mid I_{t-1}) - \E(y_t \mid I_{t-1})] - [\E(y_t \mid I_{t-2}) - \E(y_t \mid I_{t-2})] - \cdots \\
  &\quad - [\E(y_t \mid I_{t-j}) - \E(y_t \mid I_{t-j})] \\
  &= [y_t - \E(y_t \mid I_{t-1})] + [\E(y_t \mid I_{t-1}) - \E(y_t \mid I_{t-2})] + \cdots \\
  &\quad + [\E(y_t \mid I_{t-j+1}) - \E(y_t \mid I_{t-j})] + \E(y_t \mid I_{t-j}) \\
  &= (r_{t0} + r_{t1} + \cdots + r_{t,j-1}) + \E(y_t \mid y_{t-j}, y_{t-j-1}, \ldots)
\end{align*}
where
\[
  r_{tk} \equiv \E(y_t \mid I_{t-k}) - \E(y_t \mid I_{t-k-1}).
\]
This $r_{tk}$ is the revision of expectations about $y_t$ as the information set increases
from $I_{t-k-1}$ to $I_{t-k}$. By~(b), $y_t - (r_{t0} + r_{t1} + \cdots + r_{t,j-1})$ converges in mean square
to zero as $j \to \infty$ for each $t$. So $y_t$ can be written as
\[
  y_t = \sum_{j=0}^{\infty} r_{tj}.
\]
This sum is called the \textbf{telescoping sum}. The third part of Gordin's condition is
\begin{enumerate}[label=(\alph*)]
\setcounter{enumi}{2}
\item $\displaystyle \sum_{j=0}^{\infty} [\E(r_{tj}^2)]^{1/2} < \infty$.
\end{enumerate}
The telescoping sum indicates how the ``shocks'' represented by $(r_{t0}, r_{t1}, \ldots)$ influence
the current value of $y$. Condition~(c) says, roughly speaking, that shocks that
occurred a long time ago do not have a disproportionately large influence. As such,
the condition restricts the extent of serial correlation in $\{y_t\}$.

To understand the condition, take for example a zero-mean AR(1) with independent
errors:
\[
  y_t = \phi y_{t-1} + \varepsilon_t, \quad |\phi| < 1, \quad
  \{\varepsilon_t\} \text{ independent white noise with } \sigma^2 = \Var(\varepsilon_t).
\]
Evidently, Gordin's condition~(a) is satisfied. Condition~(b) too is satisfied because
\begin{equation}
  \E(y_t \mid y_{t-j}, y_{t-j-1}, \ldots) = \phi^j y_{t-j}
  \label{eq:6.5.4}
\end{equation}
and $\E[(\phi^j y_{t-j})^2] \to 0$ as $j \to \infty$. The expectation revision $r_{tj}$ can be written as
\begin{equation}
  r_{tj} = \phi^j y_{t-j} - \phi^{j+1} y_{t-j-1}
  = \phi^j (y_{t-j} - \phi y_{t-j-1}) = \phi^j \varepsilon_{t-j}.
  \label{eq:6.5.5}
\end{equation}
So the telescoping sum for AR(1) is the $\text{MA}(\infty)$ representation. Condition~(c) is
satisfied because $[\E(r_{tj}^2)]^{1/2} = |\phi|^j \sigma$ and
\[
  \sum_{j=0}^{\infty} |\phi|^j \sigma = \frac{\sigma}{1 - |\phi|} < \infty.
\]

The CLT we have been seeking is

\begin{proposition}[Gordin's CLT for zero-mean ergodic stationary processes]
\label{prop:6.10}
Suppose $\{y_t\}$ is stationary and ergodic and suppose Gordin's condition is satisfied.
Then $\E(y_t) = 0$, the autocovariances $\{\gamma_j\}$ are absolutely summable, and\footnote{This result is due to Gordin (1969), restated as Theorem~5.15 of White (1984). The claim about the absolute summability of autocovariances is noted in footnote~19 of Hansen (1982).}
\begin{equation*}
  \sqrt{n}\,\bar{y} \dto N\Bigl(0,\; \sum_{j=-\infty}^{\infty} \gamma_j\Bigr).
\end{equation*}
\end{proposition}

Since Gordin's condition is satisfied if $\{y_t\}$ is a martingale difference sequence,
which exhibits no serial correlation, Proposition~6.10 is a generalization of the
ergodic stationary Martingale Difference Sequence CLT of Chapter~2.

\subsection*{Multivariate Extension}

Extension of the above results to the multivariate case is straightforward. Let
\[
  \bar{\mathbf{y}} \equiv \frac{1}{n} \sum_{t=1}^{n} \mathbf{y}_t
\]
be the sample mean of a vector process $\{\mathbf{y}_t\}$.
\begin{itemize}
\item The multivariate version of Proposition~6.8 is that
\begin{enumerate}[label=(\alph*)]
\item $\bar{\mathbf{y}} \to_{\mathrm{m.s.}} \bm{\mu}$ if each diagonal element of $\bm{\Gamma}_j$ goes to zero as $j \to \infty$, and
\item $\lim_{n \to \infty} \Var(\sqrt{n}\,\bar{\mathbf{y}})$ (which is the \textbf{long-run covariance matrix} of $\{\mathbf{y}_t\}$) equals
$\sum_{j=-\infty}^{\infty} \bm{\Gamma}_j$ if $\{\bm{\Gamma}_j\}$ is summable (i.e., if each component of $\bm{\Gamma}_j$ is summable).
\end{enumerate}

\item Therefore, the long-run covariance matrix of $\{\mathbf{y}_t\}$ can be written as
\begin{equation}
  \mathbf{G}_Y(1) = 2\pi\,\mathbf{s}_Y(0) = \sum_{j=-\infty}^{\infty} \bm{\Gamma}_j
  = \bm{\Gamma}_0 + \sum_{j=1}^{\infty} (\bm{\Gamma}_j + \bm{\Gamma}_j'),
  \label{eq:6.5.6}
\end{equation}
where $\mathbf{G}_Y(z)$ (the autocovariance-generating function) and $\mathbf{s}_Y(\omega)$ (the spectrum)
for the vector process $\{\mathbf{y}_t\}$ are defined in Section~6.3.

\item The multivariate version of Proposition~6.9 is that
\begin{equation}
  \sqrt{n}\,(\bar{\mathbf{y}} - \bm{\mu}) \dto N\Bigl(\mathbf{0},\; \sum_{j=-\infty}^{\infty} \bm{\Gamma}_j\Bigr)
  \label{eq:6.5.7}
\end{equation}
if $\{\mathbf{y}_t\}$ is vector $\text{MA}(\infty)$ with absolutely summable coefficients and $\{\bm{\varepsilon}_t\}$ is vector
independent white noise.

\item The multivariate extension of Gordin's condition on ergodic stationary processes
is obvious:\footnote{Stated as Assumption~3.5 in Hansen (1982).}

\textbf{Gordin's condition on ergodic stationary processes}
\begin{enumerate}[label=(\alph*)]
\item $\E(\mathbf{y}_t \mathbf{y}_t')$ exists and is finite.
\item $\E(\mathbf{y}_t \mid y_{t-j}, y_{t-j-1}, \ldots) \to_{\mathrm{m.s.}} \mathbf{0}$ as $j \to \infty$.
\item $\displaystyle \sum_{j=0}^{\infty} [\E(\mathbf{r}_{tj} \mathbf{r}_{tj}')]^{1/2} < \infty$, where
\[
  \mathbf{r}_{tj} \equiv \E(\mathbf{y}_t \mid y_{t-j}, y_{t-j-1}, \ldots)
  - \E(\mathbf{y}_t \mid y_{t-j-1}, y_{t-j-2}, \ldots).
\]
\end{enumerate}

\item The multivariate version of Proposition~6.10 is obvious: Suppose Gordin's condition
holds for vector ergodic stationary process $\{\mathbf{y}_t\}$. Then $\E(\mathbf{y}_t) = \mathbf{0}$, $\{\bm{\Gamma}_j\}$ is
absolutely summable, and
\begin{equation}
  \sqrt{n}\,\bar{\mathbf{y}} \dto N\Bigl(\mathbf{0},\; \sum_{j=-\infty}^{\infty} \bm{\Gamma}_j\Bigr).
  \label{eq:6.5.8}
\end{equation}
\end{itemize}

\subsection*{Questions for Review}

\begin{enumerate}
\item Show that $\Var(\sqrt{n}\,\bar{y}) = \mathbf{1}'\,\Var(y_t, y_{t-1}, \ldots, y_{t-n+1})\,\mathbf{1}/n$, where $\Var(y_t,
y_{t-1}, \ldots, y_{t-n+1})$ is the $n \times n$ autocovariance matrix of a covariance-stationary
process $\{y_t\}$.
\end{enumerate}
